\chapter{Background and Preliminaries}
\label{chap:background}

This chapter fixes the microarchitectural and architectural background the thesis relies on,
and the model-based relational-testing framework it extends. It is written to be self-contained:
a reader who knows general computer architecture but neither the transient-execution literature
nor the AArch64 security extensions should be able to follow the contributions from here alone.
\autoref{sec:bg-speculation} and \autoref{sec:bg-sidechannels} recall out-of-order and
speculative execution and the side channels that make its transient effects observable.
\autoref{sec:bg-spectre} works through the transient-execution attacks, with a detailed
Spectre-v1 example. \autoref{sec:bg-aarch64} describes the AArch64 isolation boundaries, the
layout of a 64-bit pointer, and the security extensions---Pointer Authentication and Memory
Tagging---whose preconditions this thesis fuzzes, and \autoref{sec:bg-arm-attacks} the
hand-crafted attacks against them. \autoref{sec:bg-mrt} introduces speculation contracts,
model-based relational testing, and the \revizor{} fuzzer this work ports and generalizes.

\section{Out-of-order and speculative execution}
\label{sec:bg-speculation}

A processor core executes far faster than memory responds: a hit in the first-level cache
returns in a few cycles, while a miss that reaches main memory costs hundreds. A core that
stalled on every unresolved dependency would spend most of its time idle. High-performance cores
therefore execute \emph{out of order}---they begin any instruction whose inputs are ready rather
than waiting for its program predecessors---and retire the results \emph{in order}, so that the
architectural state changes exactly as the program specifies. The reorder buffer holds the
in-flight instructions until they can retire, and an exception or a misprediction discarded
before retirement leaves the architectural state untouched.

To keep the window of ready instructions full, the core also \emph{speculates}: it predicts the
outcome of an unresolved decision and continues past it before the decision is known. The most
important predictions are the following.
\begin{itemize}
  \item \textbf{Conditional-branch direction.} A branch predictor~\cite{tage} guesses whether a
    conditional branch is taken, so the core can keep fetching past it. The guess is learned from
    the branch's history, and an attacker who shares the predictor can train that history.
  \item \textbf{Indirect-branch and return target.} For an indirect branch or a function return,
    the core predicts the target address from a branch-target buffer (BTB) or a return-stack
    buffer, and speculatively executes at the predicted target.
  \item \textbf{Store-to-load forwarding (memory disambiguation).} When a load's address may
    alias an older, not-yet-resolved store, the core may predict no alias and forward stale
    memory to the load, or forward the store's data before the address check completes.
\end{itemize}
When a prediction is later found wrong, the core \emph{squashes} the instructions it executed on
the mistaken path---the \emph{transient} instructions---and rolls the architectural state, the
registers and memory, back to the last correct point. Program behavior is unaffected; the
transient instructions never retire.

The microarchitectural state, however, is \emph{not} rolled back. A cache line pulled in during a
transient load stays in the cache, a translation cached in a TLB stays cached, an entry allocated
in a branch predictor or a prefetcher survives, and so do effects in structures that are not
publicly documented. This asymmetry---architectural results discarded, microarchitectural effects
retained---is what every transient-execution attack exploits: a transient instruction that never
``happened'' architecturally can still leave a measurable trace.

\paragraph{A first example.}
Consider a bounds-checked array read, \reg{if (i < n) t = a[i];}, where the branch has almost
always been taken and the predictor expects it to be taken again. Call it once with \reg{i}
out of bounds. The predictor still guesses ``taken,'' so the core speculatively computes
\reg{a[i]} for the illegal index and issues the load, reading whatever lies past the array, before
the comparison \reg{i < n} resolves. When it resolves false, the load and everything after it are
squashed and \reg{t} is never written. Architecturally nothing happened. But the out-of-bounds
line \reg{a[i]} was brought into the cache and remains there, and the rest of this chapter is
about turning that residue into a readable signal.

\section{Microarchitectural side channels}
\label{sec:bg-sidechannels}

A surviving microarchitectural effect becomes a \emph{leak} once an attacker can observe it.
Observation is almost always indirect, by timing: the attacker shares a microarchitectural
structure with the victim and infers the victim's effect on that structure from the latency of
its own accesses to it. Because the shared structure is finite and its state is visible through
timing, one party's activity is measurable by the other even when the architecture forbids them
from reading each other's data.

The best-known carriers are the data caches. Two complementary techniques dominate.
\begin{description}
  \item[Prime+Probe~\cite{primeprobe}.] The attacker fills (\emph{primes}) a cache set with its
    own lines, lets the victim run, then times accesses to those lines (\emph{probe}). A line that
    is now slow was evicted by the victim, revealing that the victim touched that set. It needs no
    memory shared with the victim, only a shared cache.
  \item[Flush+Reload~\cite{flushreload}.] The attacker flushes a line it shares with the victim
    (for instance a line of a shared library or a shared read-only page), lets the victim run, and
    times a reload. A \emph{fast} reload means the victim re-fetched the line; a \emph{slow} reload
    means it did not (\autoref{fig:bg-flushreload}). It is more precise than Prime+Probe, resolving
    a single line, but requires shared memory.
\end{description}
Flush+Flush~\cite{flushflush} is a stealthier variant that times the flush itself rather than a
reload. The same timing principle applies to the translation lookaside buffers, the branch
predictors, and other shared structures: wherever two domains compete for a finite resource, the
occupancy left by one is a timing signal to the other.

\begin{figure}[t]
\centering
\fitwidth{\begin{tikzpicture}[
  font=\footnotesize,
  ph/.style={draw=black!60, thick, rounded corners=2.5pt, fill=black!3,
             align=center, minimum width=36mm, minimum height=15mm, inner sep=3pt},
  line/.style={draw=black!55, minimum width=7mm, minimum height=5mm, inner sep=0pt, font=\scriptsize},
  cached/.style={line, fill=probe!25, draw=probe!60!black},
  flushed/.style={line, fill=black!5},
  ar/.style={-{Stealth[length=2mm]}, thick, draw=black!70},
  lbl/.style={font=\scriptsize, text=black!65, align=center},
]
  % three phases, left to right
  \node[ph] (p1) at (0,0)
     {\textbf{1. Flush}\\ {\scriptsize evict the shared line}};
  \node[ph] (p2) at (4.4,0)
     {\textbf{2. Victim runs}\\ {\scriptsize accesses line if secret\,$=1$}};
  \node[ph] (p3) at (8.8,0)
     {\textbf{3. Reload + time}\\ {\scriptsize measure access latency}};
  \draw[ar] (p1) -- (p2);
  \draw[ar] (p2) -- (p3);
  % timing comparison
  \def\ty{-2.75}
  \draw[ar] (0.0,\ty) -- (10.8,\ty) node[lbl, anchor=west, text=black!70]{access latency};
  \draw[black!45, dashed] (3.1,{\ty-0.15}) -- (3.1,{\ty+1.55}) node[lbl, above, text=black!60]{threshold};
  \fill[probe!35, draw=probe!60!black] (0.1,{\ty+0.15}) rectangle (2.3,{\ty+0.65});
  \node[lbl, anchor=west, text=probe!45!black] at (2.4,{\ty+0.4}) {\textbf{fast} $\Rightarrow$ cache \emph{hit}: victim accessed it (bit $=1$)};
  \fill[culprit!22, draw=culprit!60] (0.1,{\ty+0.95}) rectangle (6.6,{\ty+1.45});
  \node[lbl, anchor=west, text=culprit] at (6.7,{\ty+1.2}) {\textbf{slow} $\Rightarrow$ cache \emph{miss}: untouched (bit $=0$)};
\end{tikzpicture}
}
\caption{Flush+Reload. The attacker evicts a line it shares with the victim, lets the victim run,
and then times a reload. A fast reload (cache hit) means the victim accessed the line; a slow
reload (miss) means it did not. Comparing the latency against a threshold recovers one secret-%
dependent bit.}
\label{fig:bg-flushreload}
\end{figure}

\paragraph{Encoding a secret into the cache.}
A single shared line carries one bit---touched or not. To carry a whole byte, an attacker uses a
\emph{probe array} of $256$ lines, one per possible value, each on its own cache line so that
touching one does not disturb the others. Code of the form \reg{t = probe[secret * 64];} brings in
exactly the line indexed by \reg{secret}; the attacker then Flush+Reloads all $256$ lines and the
one that reloads fast names the byte. This $256$-way encoding is the standard transmitter at the
end of a transient-execution attack, and the stride (here $64$ bytes, a typical line size) keeps
each value on a distinct line and defeats adjacent-line prefetching.

\paragraph{Leaks as a cascade.}
It is useful to view a leak as a cascade. A secret-dependent operation changes the state of some
\emph{carrier} structure. That change may propagate through further carriers until it reaches an
\emph{observation medium} the attacker monitors directly. The carrier and the medium need not
coincide: a victim's secret-dependent footprint may first land in the branch-target buffer and
reach the attacker-monitored L1 data cache only later, possibly in a subsequent execution. This
distinction---between where a leak is first written and where it is finally read---is central to
\autoref{chap:priming} and \autoref{chap:canon}, where the carrier is a structure other than the
monitored cache and the signal surfaces only on a later run.

\section{Transient-execution attacks}
\label{sec:bg-spectre}

Spectre~\cite{spectre} and Meltdown~\cite{meltdown} were the first attacks to combine transient
execution with a side channel to cross a security boundary. They split into two families by what
makes the offending instruction transient. \emph{Spectre}-type attacks steer a victim's own
speculation down a path that accesses secret data and transmits it through a carrier; the victim
executes the leak itself, on a mispredicted path. \emph{Meltdown}-type attacks exploit transient
execution past a fault, reading data a permission check would forbid in the brief window before
the fault is delivered.

\paragraph{A worked Spectre-v1.}
The canonical Spectre-v1 gadget (\autoref{fig:bg-spectrev1}) is a bounds-checked read feeding a
probe-array access:
\begin{quote}\ttfamily
if (x < len)\ \ y = arr2[\,arr1[x] * 64\,];
\end{quote}
The attacker first trains the branch taken by calling the gadget many times with in-bounds
\reg{x}, then flushes \reg{len} and the \reg{arr2} probe array and calls it once with \reg{x} far
out of bounds. The predictor guesses ``taken,'' so on the transient path the core reads the secret
byte $s=\reg{arr1[x]}$ from out of bounds, then reads \reg{arr2[}$s\cdot 64$\reg{]}, bringing probe
line $s$ into the cache. The comparison then resolves false and the whole path is squashed---%
architecturally \reg{y} is never written---but probe line $s$ stays cached, and a Flush+Reload
sweep of \reg{arr2} reads out $s$. Sweeping \reg{x} peels the victim's memory one byte at a time.

\begin{figure}[t]
\centering
\fitwidth{\begin{tikzpicture}[
  font=\footnotesize,
  code/.style={draw=black!45, fill=black!3, rounded corners=2pt, align=left,
               inner sep=4pt, font=\ttfamily\footnotesize},
  ev/.style={draw=black!60, thick, rounded corners=2pt, align=center,
             minimum height=9mm, inner sep=2.5pt, fill=black!4, text width=24mm},
  spec/.style={ev, draw=culprit!70, fill=culprit!8},
  cell/.style={draw=black!45, minimum width=5mm, minimum height=5mm, inner sep=0pt, font=\tiny},
  hot/.style={cell, fill=culprit!22, draw=culprit!65},
  cold/.style={cell, fill=black!3},
  ar/.style={-{Stealth[length=1.8mm]}, semithick, draw=black!70},
  sar/.style={-{Stealth[length=1.8mm]}, semithick, draw=culprit!70},
  lbl/.style={font=\scriptsize, text=black!65, align=center},
]
  % gadget code
  \node[code, anchor=north west] (gad) at (0,3.0)
     {if (x < len)\ \ \ \ \ \ // bound check\\ \ \ \ y = arr2[\,arr1[x]\,*\,64\,];};
  \node[lbl, anchor=west, text=culprit] at (0.1,1.95)
     {branch \emph{mistrained} taken; call it with out-of-bounds \texttt{x}};
  % transient timeline
  \node[spec] (e1) at (1.4,0.7)  {predict\\ \texttt{x<len} taken};
  \node[spec] (e2) at (4.4,0.7)  {load \texttt{arr1[x]}\\ $\to$ secret $s$};
  \node[spec] (e3) at (7.4,0.7)  {load \texttt{arr2[$s{\cdot}64$]}\\ brings in line $s$};
  \node[ev]   (e4) at (10.4,0.7) {resolve: \texttt{x$\ge$len}\\ \emph{squash}};
  \draw[sar] (e1) -- (e2);
  \draw[sar] (e2) -- (e3);
  \draw[ar]  (e3) -- (e4);
  \node[lbl, anchor=north, text=culprit] at (4.4,0.05) {transient path (rolled back architecturally)};
  % cache footprint
  \node[lbl, anchor=east] at (-0.15,-1.75) {\texttt{arr2} lines:};
  \foreach \i in {0,...,11}{
     \ifnum\i=7 \node[hot] (c\i) at ({0.3+\i*0.56},-1.75) {$s$};
     \else \node[cold] (c\i) at ({0.3+\i*0.56},-1.75) {\i};\fi}
  \draw[sar] (e3.south) to[out=-90,in=90] (c7.north);
  \node[lbl, anchor=west, text=culprit] at ({0.3+12*0.56},-1.75)
     {attacker Flush+Reloads \texttt{arr2}: the hot line reveals $s$};
\end{tikzpicture}
}
\caption{A worked Spectre-v1. A mistrained bounds check leaves a secret-dependent load on the
transient path. The transient \reg{arr2} access caches a line indexed by the secret byte $s$;
after the misprediction is squashed, the attacker recovers $s$ by Flush+Reloading \reg{arr2}.
Architecturally the two calls that differ only in the secret are identical---the signal lives
entirely in the surviving cache state.}
\label{fig:bg-spectrev1}
\end{figure}

The example previews the relational view of \autoref{sec:bg-mrt}: two victim runs that differ only
in a secret the program never reveals architecturally---here two values of the out-of-bounds
byte---produce identical architectural behavior yet different cache footprints. That gap between
``architecturally indistinguishable'' and ``distinguishable through the microarchitecture'' is
exactly what a contract violation will formalize.

\paragraph{The Spectre variants.}
The family is distinguished by the speculation primitive it abuses.
\begin{itemize}
  \item \textbf{Spectre-v1} (bounds-check bypass) mistrains a conditional branch so a bounds check
    is bypassed on the transient path, as above.
  \item \textbf{Spectre-v2} (branch-target injection) poisons an indirect branch's predicted
    target so the victim speculatively executes an attacker-chosen gadget. Two sub-cases matter
    for this thesis: the gadget's transient memory accesses leave a footprint in the \emph{L1 data
    cache}, the channel standard tooling already probes; or the poisoned prediction itself leaves
    its mark only in the \emph{branch-target buffer}, a channel that surfaces only on re-execution
    (\autoref{chap:priming}, \autoref{chap:canon}).
  \item \textbf{Spectre-v4} (speculative store bypass) exploits memory disambiguation: a load is
    speculatively forwarded stale data because the aliasing store has not yet resolved, running a
    dependent access on the transient path.
  \item \textbf{Spectre-RSB} drives speculation through the return-address stack, running a return
    along a stale predicted target.
\end{itemize}
\autoref{chap:canon} uses Spectre-v1, Spectre-v4, and Spectre-RSB interchangeably as windows that
reach the same precondition check transiently, which is what makes the leak it studies robust to
the particular primitive a target exposes.

\section{The AArch64 security architecture}
\label{sec:bg-aarch64}

AArch64 is the 64-bit execution state of the Arm architecture. This section covers the parts the
thesis depends on: the privilege boundaries, the shared microarchitecture across them, the layout
of a virtual address, and the two security extensions whose operand checks the precondition
fuzzing of \autoref{chap:precond} targets.

\paragraph{Exception levels and translation regimes.}
AArch64 separates software by privilege into \emph{exception levels}: EL0 for user applications,
EL1 for operating-system kernels, EL2 for hypervisors, and EL3 for secure
firmware~\cite{arm_exception_levels} (\autoref{fig:bg-elmap}). A virtual address is translated
through one of two regimes, selected by its top bits: the low range translates through \reg{TTBR0}
and the high range through \reg{TTBR1}~\cite{arm64pagetables}. These boundaries isolate
architectural state, but the microarchitecture is only partially isolated. TLB entries are
typically tagged with an address-space identifier (ASID) and, under virtualization, a
virtual-machine identifier (VMID), so one domain cannot \emph{use} another's entries. The branch
predictors are not specified by the architecture~\cite{chen2024dissectingconditionalbranchpredictors,jia_cbp},
though cores implementing FEAT\_CSV2 prevent branch targets trained in one context from steering
speculation in another. Tagging and context isolation prevent one domain from using another's
state, but unless a structure is partitioned the domains still compete for its capacity, so
activity in one domain can change state that another domain can measure.

\begin{figure}[t]
\centering
\fitwidth{\begin{tikzpicture}[
  font=\footnotesize,
  el/.style={draw=black!65, thick, rounded corners=2pt, align=center,
             minimum width=36mm, minimum height=9mm, inner sep=2pt},
  reg/.style={draw=black!55, thick, align=center, inner sep=2pt, minimum width=38mm},
  lbl/.style={font=\scriptsize, text=black!65, align=center},
  ar/.style={-{Stealth[length=2mm]}, thick, draw=black!70},
]
  % left: exception-level stack (EL0 bottom .. EL3 top)
  \node[el, fill=probe!14]  (el0) at (0,0)     {\textbf{EL0}\ \ {\scriptsize user applications}};
  \node[el, fill=laneB!16]  (el1) at (0,1.1)   {\textbf{EL1}\ \ {\scriptsize OS kernel}};
  \node[el, fill=laneA!16]  (el2) at (0,2.2)   {\textbf{EL2}\ \ {\scriptsize hypervisor (VHE)}};
  \node[el, fill=culprit!12](el3) at (0,3.3)   {\textbf{EL3}\ \ {\scriptsize secure firmware}};
  \draw[ar] (-2.45,-0.45) -- (-2.45,3.75)
       node[lbl, rotate=90, anchor=south, pos=0.5, yshift=4pt]{increasing privilege};

  % right: 64-bit address space, the two VA ranges
  \def\rx{6.0}
  \node[reg, fill=laneB!16, minimum height=11mm, anchor=south] (hi) at (\rx,2.4)
     {\textbf{\texttt{TTBR1}} range (high)\\{\scriptsize top bits all \textbf{1}, kernel}};
  \node[reg, fill=black!8, minimum height=12mm, anchor=south] (hole) at (\rx,0.7)
     {\emph{non-canonical hole}\\{\scriptsize any other top bits, faults}};
  \node[reg, fill=probe!14, minimum height=11mm, anchor=south] (lo) at (\rx,-1.0)
     {\textbf{\texttt{TTBR0}} range (low)\\{\scriptsize top bits all \textbf{0}, user}};
  \node[lbl, anchor=west] at ({\rx+2.0},3.3)  {\texttt{0xFFFF\_FFFF\_FFFF\_FFFF}};
  \node[lbl, anchor=west] at ({\rx+2.0},-0.95){\texttt{0x0000\_0000\_0000\_0000}};
  \draw[ar] (el0.east) to[bend right=10]
       node[lbl, fill=white, inner sep=1.5pt, pos=0.5]{issues VA} (lo.west);
  \draw[ar] (el1.east) to[bend left=10]
       node[lbl, fill=white, inner sep=1.5pt, pos=0.5]{issues VA} (hi.west);

  % bottom: how the microarchitecture is (only partially) isolated across privilege
  \node[draw=black!45, rounded corners=3pt, fill=black!2, inner sep=6pt,
        anchor=north] (tbl) at (3.0,-1.95) {%
    \begin{tabular}{@{}l@{\hspace{7mm}}l@{}}
    \multicolumn{2}{@{}l@{}}{\textbf{\footnotesize microarchitectural isolation across privilege}}\\[3pt]
    caches (L1, L2) & shared (address tag, not per-domain) \\
    TLB            & tagged per domain (\reg{ASID}, \reg{VMID}) \\
    BTB            & physically split \\
    \end{tabular}};
\end{tikzpicture}
}
\caption{AArch64 exception levels (left) and the two translation regimes (right). Architectural
state is isolated per level; microarchitectural structures are shared or merely tagged. The top
address bits select the regime, and canonicality requires those bits to be uniformly zero
(\reg{TTBR0}) or one (\reg{TTBR1}); any other value is out of range and faults.}
\label{fig:bg-elmap}
\end{figure}

\subsection{Virtual addressing, canonicality, and the top byte}
\label{sec:bg-canon}

An address is \emph{canonical} when the bits above the configured input-address size are all zero,
selecting the \reg{TTBR0} range, or all one, selecting the \reg{TTBR1} range. Any other value is
out of range and raises an address-size or translation fault. Arm does not use the x86 term
``canonical'' for this condition, but the requirement is its analogue and we use the term
throughout. A 64-bit pointer therefore carries meaning in bits well above the virtual-address
field, and the Arm extensions place further data there. \autoref{fig:bg-ptrlayout} shows the
layout for a 48-bit virtual address.

\begin{figure}[t]
\centering
\fitwidth{\begin{tikzpicture}[
  font=\footnotesize,
  bit/.style={draw=black!65, thick, minimum height=8mm, inner sep=0pt, align=center},
  va/.style={bit, fill=probe!18, draw=probe!55!black},
  tbi/.style={bit, fill=black!8},
  mte/.style={bit, fill=laneB!22, draw=laneB!60!black},
  reg/.style={bit, fill=culprit!16, draw=culprit!70},
  pac/.style={bit, fill=laneA!20, draw=laneA!65!black},
  br/.style={decorate, decoration={brace, amplitude=4pt, raise=2pt}, draw=black!60},
  tk/.style={font=\scriptsize, text=black!65},
  lbl/.style={font=\scriptsize, align=center},
]
  \def\u{0.52}            % cm per bit (upper 16 bits, to scale)
  \def\h{0.8}             % box height in cm
  % upper 16 bits [63:48], MSB on the left
  % MTE/top byte: bits 63..56 => x in [0, 8u]; split [63:60] and [59:56]
  \node[tbi, minimum width={4*\u cm}, anchor=south west] (b6360) at (0,0) {\scriptsize$63{:}60$};
  \node[mte, minimum width={4*\u cm}, anchor=south west] (b5956) at ({4*\u},0) {\scriptsize$59{:}56$};
  \node[reg, minimum width={1*\u cm}, anchor=south west] (b55)   at ({8*\u},0) {\scriptsize$55$};
  \node[pac, minimum width={7*\u cm}, anchor=south west] (b5448) at ({9*\u},0) {\scriptsize$54{:}48$};
  % break marker (two short slashes)
  \draw[black!45,line width=0.8pt] (8.46,-0.08) -- (8.62,0.88);
  \draw[black!45,line width=0.8pt] (8.60,-0.08) -- (8.76,0.88);
  % VA [47:0], not to scale
  \node[va, minimum width=5.6cm, minimum height={\h cm}, anchor=south west] (bva) at ({16*\u+0.46},0)
     {bits $[47{:}0]$, the virtual address (48-bit VA)};
  % field braces/labels above
  \draw[br] (b6360.north west) -- (b5956.north east)
       node[midway, above=5pt, lbl, text=black!70]{top byte $[63{:}56]$\\{\scriptsize ignored by translation when \texttt{TBI}=1}};
  \draw[br] (b5956.north west) -- (b5956.north east)
       node[midway, above=20pt, lbl, text=laneB!55!black]{\reg{MTE} logical\\ tag $[59{:}56]$};
  % labels below
  \node[lbl, text=culprit, anchor=north, align=center] at ($(b55.south)+(0,-0.12)$)
     {bit $55$: regime select\\(\texttt{TTBR0}/\texttt{TTBR1})\\ effective sign bit under \texttt{TBI}};
  \draw[br, decoration={mirror}] ($(b6360.south west)+(0,-1.55)$) -- ($(b5448.south east)+(0,-1.55)$) ;
  \node[lbl, text=laneA!55!black, anchor=north, align=center] at ({8*\u},-2.35)
     {\reg{PAC} field: the spare upper bits, excluding bit $55$\\{\scriptsize bits $[54{:}48]$, plus the top byte $[63{:}56]$ when \texttt{TBI} is off}};
  % bit ruler ticks
  \foreach \b/\x in {63/0, 56/{8*\u}, 55/{9*\u}, 48/{16*\u}}
     \node[tk, anchor=south] at ({\x},{\h+0.02}) {};
\end{tikzpicture}
}
\caption{Layout of a 64-bit AArch64 pointer with a 48-bit virtual address. The low bits hold the
address; the upper bits are shared among canonicality, the ignored top byte, the MTE logical tag,
and the PAC field. Bit~55 selects the translation regime and is the effective sign bit when Top
Byte Ignore is enabled (bit~63 when it is not); it is excluded from the PAC field.}
\label{fig:bg-ptrlayout}
\end{figure}

\paragraph{Top Byte Ignore (TBI).}
When TBI is enabled, the top byte of a pointer, bits [63:56], is excluded from translation and
does not participate in the canonicality check, so software may store metadata there and the
hardware strips it before translating~\cite{arm_tbi}. The effective sign bit that must be
replicated to stay canonical then drops to bit~55; with TBI disabled the top byte participates and
the sign bit is bit~63. This TBI-dependence of the effective sign bit recurs in
\autoref{chap:precond} and \autoref{chap:canon}, where a primitive must flip the regime-correct
bit.

\paragraph{Pointer Authentication (PAuth).}
Introduced in Armv8.3-A, PAuth protects pointer integrity by storing a pointer authentication code
(PAC) in a pointer's otherwise unused upper bits, excluding bit~55, which selects the translation
regime~\cite{arm_pac,pac_whitepaper}. The PAC width therefore depends on the virtual-address size
and the TBI configuration (\autoref{fig:bg-ptrlayout}). The code is computed from the pointer, a
64-bit modifier, and a secret key, using the QARMA block cipher. The \reg{PAC*} instructions
compute and insert it; the \reg{AUT*} instructions verify it, and on success restore the upper bits
to the sign extension of bit~55. Combined instructions such as \reg{RETAA} authenticate and branch
in one step. A failed authentication yields a corrupted, non-canonical pointer that faults on use,
or faults immediately on cores implementing FEAT\_FPAC and FEAT\_FPACCOMBINE.

\paragraph{Memory Tagging Extension (MTE).}
Introduced in Armv8.5-A, MTE associates a 4-bit \emph{allocation tag} with every 16-byte tag
granule of physical memory~\cite{arm_mte,mte_whitepaper}, while a pointer carries a 4-bit
\emph{logical tag} in bits [59:56], inside the top byte that TBI excludes from translation
(\autoref{fig:bg-ptrlayout}). The allocation tags live in a separate tag memory, accessed with
dedicated instructions such as \reg{LDG} and \reg{STG}. On each access to tagged normal memory the
processor compares the pointer's logical tag against the memory's allocation tag; a mismatch raises
a synchronous or asynchronous tag-check fault, or is ignored, depending on configuration. MTE is a
probabilistic defense against spatial and temporal memory-safety errors: an out-of-bounds or
use-after-free access is likely to carry the wrong tag and be caught.

\section{Arm-specific speculative attacks}
\label{sec:bg-arm-attacks}

Two published attacks turn these extensions against themselves by checking their preconditions on a
transient path, where a failed check does not architecturally fault.

PACMAN~\cite{pacman} bypasses PAuth by testing PAC guesses during transient execution. A correct
guess authenticates and the pointer is used speculatively, leaving a microarchitectural footprint;
an incorrect guess fails authentication, which does not fault on the transient path but does
suppress the footprint. The presence or absence of the footprint is thus an oracle on whether a PAC
guess was correct, which removes the architectural crash that would otherwise make brute-forcing a
PAC infeasible.

TikTag~\cite{tiktag} does the analogous thing for MTE: it leaks allocation tags by observing
tag-check outcomes during transient execution, where a mismatch does not architecturally fault. A
speculative access with a matching logical tag proceeds and leaves a footprint; a mismatch is
suppressed, and the footprint difference reveals the granule's allocation tag, defeating the
randomness MTE relies on.

PACMAN and TikTag share a structure this thesis generalizes. Each targets a \emph{precondition}
that an instruction places on its operand---a valid PAC, a matching MTE tag---and each observes,
through a side channel, the outcome of the precondition check on a transient path where its
violation does not fault. Both were built by hand, for one precondition each. The methodology of
\autoref{chap:precond} fuzzes preconditions of this kind systematically, and \autoref{chap:canon}
applies it to a precondition---address canonicality and alignment---that is not a defensive
extension at all but a basic property of address translation.

\section{Speculation contracts and model-based relational testing}
\label{sec:bg-mrt}

The attacks above are found and understood by hand. To search for leaks automatically we need a
precise, implementation-independent statement of what a CPU is \emph{allowed} to leak, and a way
to test a black-box CPU against it. Speculation contracts provide the specification, and
model-based relational testing provides the search.

\paragraph{Speculation contracts.}
A speculation contract~\cite{contracts} specifies what a CPU may leak without modeling its
implementation. It has an \emph{execution clause}, specifying which instructions may trigger
speculation and how (for example, that conditional branches may be mispredicted), and an
\emph{observation clause}, specifying what each executed instruction exposes (for example, the
addresses of its loads and stores)~\cite{revizor}. Executing a program $P$ on an input $I$ under a
contract $\mathcal{C}$ yields a \emph{contract trace} $\mathit{CTrace}=\mathcal{C}(P,I)$. The
baseline contract used throughout this thesis is \textsc{CT-SEQ}, under which $\mathcal{C}(P,I)$ is
the sequence of memory addresses and control-flow targets of the in-order, non-speculative
execution. Intuitively, \textsc{CT-SEQ} permits a constant-time program's ordinary address trace
and forbids any leak that appears only under speculation.

\paragraph{Hardware traces.}
On the CPU, the same execution starts from a microarchitectural context $\mathit{Ctx}$---the
initial state of the caches, buffers, and predictors---and passes through a sequence of
microarchitectural states $\mu_0(=\mathit{Ctx}),\mu_1,\dots,\mu_n$. An attacker cannot read these
states; it observes them only through a side channel, modeled as an observation function $O$ that
projects a state onto its attacker-visible part~\cite{contracts}. The \emph{hardware trace} is the
sequence of these observations~\cite{revizor}:
\begin{equation}
\mathit{HTrace}=\mathit{Attack}_O(P,I,\mathit{Ctx}) = O(\mu_0)\cdot O(\mu_1)\cdots O(\mu_n).
\label{eq:bg-htrace}
\end{equation}
A CPU \emph{satisfies} contract $\mathcal{C}$ with respect to observer $O$ when, for every program
$P$, every pair of inputs $I_1,I_2$, and every context $\mathit{Ctx}$,
\begin{equation}
\mathcal{C}(P,I_1) = \mathcal{C}(P,I_2) \;\Longrightarrow\;
\mathit{Attack}_O(P,I_1,\mathit{Ctx}) = \mathit{Attack}_O(P,I_2,\mathit{Ctx}).
\label{eq:bg-contract}
\end{equation}
Inputs the contract deems indistinguishable must stay indistinguishable to the attacker from any
common starting state. A \emph{contract violation} is a tuple $(P,I_1,I_2,\mathit{Ctx})$ for which
\begin{equation}
\mathcal{C}(P,I_1) = \mathcal{C}(P,I_2) \;\land\;
\mathit{Attack}_O(P,I_1,\mathit{Ctx}) \neq \mathit{Attack}_O(P,I_2,\mathit{Ctx}).
\label{eq:bg-violation}
\end{equation}
The worked Spectre-v1 of \autoref{sec:bg-spectre} is exactly such a tuple: the two secret-%
dependent calls share a \textsc{CT-SEQ} trace (neither performs the out-of-bounds access
architecturally) yet leave different cache footprints, so $\mathit{HTrace}$ differs where
$\mathit{CTrace}$ does not.

\paragraph{Model-based relational testing.}
Model-based relational testing (MRT) searches for violations by sampling. It generates a program
$P$ and a set of inputs, computes contract traces on a model of $\mathcal{C}$ to partition the
inputs into \emph{contract-equivalence classes}, measures hardware traces on the CPU under test,
and reports a \emph{candidate violation} when two inputs of one class produce different hardware
traces. A differing trace is not by itself evidence of a leak:
Equation~\eqref{eq:bg-violation} requires both inputs to run from the same context, and on a real
CPU the context cannot be set directly, so each candidate must be validated to show the difference
is caused by the inputs rather than by their differing contexts.

\paragraph{The Revizor pipeline.}
\revizor{}~\cite{revizor,revizor_micro} implements MRT as four components followed by a validation
step (\autoref{fig:bg-revizor}). A \emph{generator} produces small random programs, with memory
accesses confined to a sandboxed region, together with random inputs that initialize the registers
and the sandbox memory. A \emph{model} runs each program and input on an emulator of the contract
and records the contract trace. An \emph{executor} runs them on the CPU under test and records the
hardware trace. An \emph{analyzer} groups inputs into contract-equivalence classes and reports any
within-class hardware-trace mismatch as a candidate. A validation step, \emph{priming}, then
re-executes inputs to decide whether each candidate is a true violation or a false positive.

\begin{figure}[t]
\centering
\fitwidth{\begin{tikzpicture}[
  font=\footnotesize,
  comp/.style={draw=black!70, fill=black!4, rounded corners=2.5pt, thick,
               align=center, minimum width=20mm, minimum height=7mm, inner sep=1.5pt},
  stage/.style={draw=black!70, dashed, fill=white, rounded corners=2.5pt, thick,
                align=center, minimum width=20mm, minimum height=7mm, inner sep=1.5pt},
  out/.style={text=#1, font=\footnotesize\itshape, inner sep=1.5pt},
  lbl/.style={font=\scriptsize, text=black!80, inner sep=1.5pt},
  arr/.style={-{Stealth[length=1.8mm,width=1.5mm]}, thick, draw=black!75},
]
  \node[comp] (gen)   at (0, 0)       {Generator};
  \node[comp] (model) at (-1.7,-1.45) {Model\\[-2pt]{\scriptsize emulator of $\mathcal{C}$}};
  \node[comp] (exec)  at ( 1.7,-1.45) {Executor\\[-2pt]{\scriptsize CUT, $O_{\mathrm{L1D}}$}};
  \node[comp]  (an)    at (0,-2.95) {Analyzer};
  \node[stage] (prime) at (0,-4.2)  {Priming\\[-2pt]{\scriptsize validation step}};
  \node[out=red!75!black] (viol) at (-1.3,-5.5) {violation};
  \node[out=black!65]    (fp)   at ( 1.3,-5.5) {false positive};
  \coordinate (split) at (0,-0.8);
  \draw[thick, draw=black!75] (gen.south) -- (split);
  \draw[arr] (split) -| (model.north);
  \draw[arr] (split) -| (exec.north);
  \node[lbl, right] at ($(gen.south)!0.5!(split)$) {$P,\, I_1,\dots,I_N$};
  \draw[arr] (model.south) |- (an.west);
  \draw[arr] (exec.south)  |- (an.east);
  \node[lbl, left, align=right] at ($(model.south)!0.4!(model.south |- an.west)$) {contract\\traces};
  \node[lbl, right, align=left] at ($(exec.south)!0.4!(exec.south |- an.east)$)   {hardware\\traces};
  \draw[arr] (an) -- node[lbl, right] {candidates} (prime);
  \coordinate (split2) at (0,-4.95);
  \draw[thick, draw=black!75] (prime.south) -- (split2);
  \draw[arr] (split2) -| (viol.north);
  \draw[arr] (split2) -| (fp.north);
  \path let \p1 = (current bounding box.north east),
            \p2 = (current bounding box.south west),
            \n1 = {max(\x1, -\x2)} in
        coordinate (bbne) at (\n1, \y1) coordinate (bbsw) at (-\n1, \y2);
  \pgfresetboundingbox
  \useasboundingbox (bbsw) rectangle (bbne);
\end{tikzpicture}
}
\caption{\revizor{}'s testing pipeline: four components (solid) and the priming validation step
(dashed). The model and the executor run the same programs and inputs; the analyzer compares their
traces class by class.}
\label{fig:bg-revizor}
\end{figure}

\paragraph{Hardware traces, context, and priming.}
Measuring the full hardware trace of Equation~\eqref{eq:bg-htrace} on a black-box CPU is
infeasible, so \revizor{} measures a single structure, the L1 data cache, at the end of each
input's execution, through a cache side channel such as Flush+Reload or Prime+Probe. It therefore
tests the contract with respect to an \emph{L1D observer}: effects in other structures appear in
the hardware trace only once they propagate into the L1D. Because the context of a black-box CPU
cannot be set directly, \revizor{} sets it indirectly, running the inputs back to back with only a
partial reset between them~\cite{revizor,hide_seek}, so each input's context is the state its
predecessors left. A candidate violation is then a pair of contract-equivalent inputs, at
positions $i<j$, whose L1D traces differ. Since the two traces are measured in different contexts,
the difference may stem from the contexts rather than the inputs. \emph{Priming} closes this gap by
re-running the sequence with the two inputs swapped into each other's position, so each input is
measured in both contexts. The candidate is a false positive when both swapped comparisons agree,
and a violation otherwise. \autoref{chap:priming} shows that this rule, while sound, is incomplete,
and replaces it with a provably stronger one.

\paragraph{Boosting.}
Two inputs are \emph{contract-equivalent}, or in the same \emph{class}, for a program $P$ when they
share a contract trace. To \emph{boost} an input $I$, \revizor{} generates further inputs of the
same class---distinct inputs with an identical contract trace~\cite{revizor,hide_seek}---by running
the contract model on $I$, keeping the input components the trace depends on, and re-randomizing the
rest. Boosting populates the equivalence classes that both priming and its generalization
(\autoref{chap:priming}) compare across.
